\documentclass[11pt, a4paper]{article}

\usepackage[utf8]{inputenc}
\usepackage[T1]{fontenc}
\usepackage{amsmath, amssymb, amsthm, mathtools}
\usepackage{bm}
\usepackage{booktabs}
\usepackage{graphicx}
\usepackage{hyperref}
\usepackage[margin=2.5cm]{geometry}
\usepackage{enumitem}
\usepackage{xcolor}
\usepackage{tcolorbox}
\usepackage{float}

% Theorem environments
\theoremstyle{definition}
\newtheorem{definition}{Definition}[section]
\newtheorem{example}{Example}[section]
\theoremstyle{plain}
\newtheorem{theorem}{Theorem}[section]
\newtheorem{proposition}{Proposition}[section]
\newtheorem{corollary}{Corollary}[section]
\theoremstyle{remark}
\newtheorem{remark}{Remark}[section]

% Colored boxes
\newtcolorbox{keyresult}[1][]{
    colback=blue!5!white,
    colframe=blue!60!black,
    fonttitle=\bfseries,
    title={Key Insight},
    #1
}
\newtcolorbox{warningbox}[1][]{
    colback=red!5!white,
    colframe=red!60!black,
    fonttitle=\bfseries,
    title={Warning},
    #1
}
\newtcolorbox{refbox}[1][]{
    colback=gray!5!white,
    colframe=gray!50!black,
    fonttitle=\bfseries,
    title={Go Deeper},
    #1
}
\newtcolorbox{keyidea}[1][]{
    colback=green!5!white,
    colframe=green!50!black,
    fonttitle=\bfseries,
    title={Core Idea},
    #1
}
\newtcolorbox{prereq}[1][]{
    colback=orange!5!white,
    colframe=orange!60!black,
    fonttitle=\bfseries,
    title={Read Before Coding},
    #1
}
\newtcolorbox{objective}[1][]{
    colback=blue!5!white,
    colframe=blue!60!black,
    fonttitle=\bfseries,
    title={What You Will Build},
    #1
}

% Custom commands
\newcommand{\VaR}{\operatorname{VaR}}
\newcommand{\ES}{\operatorname{ES}}
\newcommand{\CVaR}{\operatorname{CVaR}}
\newcommand{\EV}{\mathbb{E}}
\newcommand{\PV}{\mathbb{P}}
\newcommand{\RV}{\mathbb{R}}
\newcommand{\ind}{\boldsymbol{1}}
\newcommand{\PnL}{\operatorname{PnL}}

\title{%
    \textbf{Deep Hedging} \\[0.5em]
    \large Neural Networks for Optimal Derivative Hedging
}
\author{Andr\'es Gonz\'alez Ortega}
\date{March 2026}

\begin{document}
\maketitle

\begin{abstract}
An applied treatment of deep hedging --- the use of neural networks to learn
optimal derivative hedging strategies directly from simulated data.
We begin with the classical Black--Scholes delta hedging framework, expose its failures
under realistic market conditions (discrete rebalancing, transaction costs, model misspecification),
and develop the deep hedging framework of Buehler et al.\ (2019) as a principled alternative.
A worked numerical example runs throughout: hedging a European call option over 30 days
under transaction costs, comparing Black--Scholes delta hedging against neural-network policies
trained with variance and CVaR objectives.
We assume familiarity with option pricing, stochastic calculus at the level of
Black--Scholes, and basic neural network training.
\end{abstract}

\tableofcontents
\newpage

%% ============================================================================
\section{The Hedging Problem}
\label{sec:hedging-problem}
%% ============================================================================

You sell a European call option with strike $K$ and maturity $T$ on a stock currently
trading at $S_0$, and you receive the option premium $C_0$.
Your liability at expiry is $\max(S_T - K, 0)$.
To manage this risk, you dynamically trade the underlying stock, holding
$\delta_t$ shares at each time step $t$.

Under the Black--Scholes model, the optimal hedge is to hold $\delta_t = \Delta_t^{\text{BS}}$
shares at each instant, where the Black--Scholes delta is:
\[
    \Delta_t^{\text{BS}} = \frac{\partial C}{\partial S}\bigg|_{S=S_t}
    = \Phi(d_1), \quad
    d_1 = \frac{\ln(S_t / K) + (r + \sigma^2/2)(T - t)}{\sigma\sqrt{T - t}}
\]
with $\Phi$ the standard normal CDF, $r$ the risk-free rate, and $\sigma$ the volatility.

This is optimal \textbf{if and only if} three conditions hold simultaneously:
\begin{enumerate}
    \item Volatility $\sigma$ is \emph{constant} and known.
    \item There are \emph{no transaction costs} --- rebalancing is free.
    \item Rebalancing occurs \emph{continuously} --- at every instant, not just daily.
\end{enumerate}

None of these hold in practice. The consequence: even with perfect knowledge of the
model parameters, delta hedging produces a \textbf{random} profit-and-loss (P\&L),
not a riskless one.

\subsection{The Hedging P\&L}

Partition $[0, T]$ into $N$ equally spaced rebalancing dates $t_0 = 0, t_1, \ldots, t_N = T$
with $\Delta t = T/N$.
At each $t_k$, the hedger holds $\delta_k$ shares.
The hedging P\&L at maturity is:
\[
    \PnL = C_0 - \max(S_T - K, 0) + \sum_{k=0}^{N-1} \delta_k (S_{t_{k+1}} - S_{t_k})
    - \sum_{k=0}^{N-1} \text{cost}_k
\]
where the first term is the premium received, the second is the option payoff owed,
the third is the cumulative hedging gains from trading the stock,
and the last term captures transaction costs (which we set to zero for now).

Under perfect Black--Scholes conditions with continuous rebalancing ($N \to \infty$),
$\PnL = 0$ almost surely.
With discrete rebalancing, $\PnL$ is a random variable with mean near zero
but nonzero variance.

\subsection{Worked Example: Daily Delta Hedging}

\begin{example}[Hedging a European call with daily rebalancing]\label{ex:daily-hedge}
Consider the following setup:
\begin{itemize}
    \item $S_0 = 100$, $K = 100$ (at-the-money), $r = 0$, $\sigma = 20\%$, $T = 30/365$ years.
    \item Black--Scholes price: $C_0 = \$2.18$.
    \item Rebalancing: daily ($N = 30$ steps, $\Delta t = 1/365$).
    \item The true dynamics follow GBM with $\sigma = 20\%$ (no model misspecification).
    \item Transaction costs: zero (for now).
\end{itemize}

\textbf{Simulation}: generate 10{,}000 GBM paths.
At each step, compute $\Delta_k^{\text{BS}}$ and hold that many shares.
At maturity, compute the P\&L for each path.

\textbf{Results} (representative):
\begin{center}
\begin{tabular}{@{}lc@{}}
\toprule
\textbf{Statistic} & \textbf{Hedging P\&L} \\
\midrule
Mean      & $\$0.002$ \\
Std.\ dev.\ & $\$0.41$ \\
Skewness  & $-0.35$ \\
5\% CVaR  & $-\$0.89$ \\
\bottomrule
\end{tabular}
\end{center}

The mean is near zero (good), but the standard deviation is \$0.41 --- about 19\%
of the option premium.
The negative skewness means losses are more severe than gains.
The 5\% CVaR says that in the worst 5\% of outcomes, the average loss is \$0.89,
or 41\% of the premium received.

This is \emph{without} transaction costs and \emph{without} model error.
In practice, the situation is much worse.
\end{example}

\begin{keyresult}
Even under the best possible conditions (correct model, zero costs), daily delta hedging
produces hedging P\&L with significant variance.
The hedging error scales as $O(\sqrt{\Delta t})$: doubling the rebalancing frequency
reduces the standard deviation by only $\sqrt{2} \approx 1.41$.
To cut the hedging error in half, you need to rebalance four times as often.
\end{keyresult}

%% ============================================================================
\section{Why Black--Scholes Fails in Practice}
\label{sec:bs-failures}
%% ============================================================================

The worked example above already shows that discrete rebalancing introduces hedging error.
But this is only the first of several compounding failures.

\subsection{Discrete Rebalancing}

Under continuous rebalancing, the self-financing hedging portfolio perfectly replicates
the option payoff.
With discrete steps of size $\Delta t$, each interval introduces a hedging error
proportional to the \emph{gamma} (second derivative) of the option:
\[
    \epsilon_k \approx \frac{1}{2}\Gamma_k (\Delta S_k)^2 - \frac{1}{2}\Gamma_k \sigma^2 S_k^2 \Delta t
\]
where $\Gamma_k = \partial^2 C / \partial S^2$.
This is the difference between the realized and expected quadratic variation
over each interval.
The standard deviation of the total hedging error scales as $O(\sigma^2 \sqrt{\Delta t})$.

\begin{warningbox}
Gamma is largest for at-the-money options near expiry.
This is precisely when hedging errors are most dangerous:
the delta swings rapidly between 0 and 1,
and each discrete rebalancing step misses the rapid changes.
\end{warningbox}

\subsection{Transaction Costs}

Each rebalance incurs a cost.
A simple proportional cost model charges:
\[
    \text{cost}_k = c \cdot |\delta_{k} - \delta_{k-1}| \cdot S_{t_k}
\]
where $c$ is the cost rate (e.g., $c = 0.001$ for 10 basis points).
The total transaction cost over the life of the hedge is:
\[
    \text{TC} = c \sum_{k=1}^{N} |\delta_k - \delta_{k-1}| \cdot S_{t_k}
\]

This creates a dilemma:
\begin{itemize}
    \item Rebalance \emph{more often} to reduce hedging error $\Rightarrow$ pay more in transaction costs.
    \item Rebalance \emph{less often} to save on costs $\Rightarrow$ suffer larger hedging error.
\end{itemize}

The optimal trade-off depends on the cost rate, the option's gamma profile, and
the hedger's risk preferences --- there is no universal closed-form solution.

\subsection{Model Misspecification}

If the true volatility $\sigma_{\text{true}}$ differs from the implied volatility
$\sigma_{\text{BS}}$ used to compute deltas, the hedge is systematically biased.
Specifically, if $\sigma_{\text{true}} > \sigma_{\text{BS}}$, the hedger underestimates delta
for OTM calls and overestimates it for ITM calls, leading to a systematic P\&L drag.

More generally, the Black--Scholes model assumes:
\begin{itemize}
    \item Constant volatility (reality: stochastic, with leverage effects and clustering).
    \item Continuous paths (reality: jumps occur, especially during crises).
    \item Log-normal returns (reality: heavy tails, skewness).
\end{itemize}

Under stochastic volatility (e.g., Heston), the correct hedge requires not just
the delta but also the \emph{vega hedge} --- trading another option to neutralize
volatility exposure.
Under jumps, perfect hedging is impossible even in theory.

\subsection{Exotic Options}

For exotic options --- barrier options, Asian options, lookback options --- the Greeks
may not exist in closed form.
Computing deltas requires either:
\begin{itemize}
    \item Finite-difference bumping (numerically expensive, noisy for path-dependent payoffs).
    \item Monte Carlo estimation of sensitivities (variance reduction needed).
    \item Adjoint algorithmic differentiation (efficient but complex to implement).
\end{itemize}

In all these cases, even obtaining the hedging rule is expensive,
let alone optimizing it for transaction costs and risk preferences.

\begin{keyresult}
The practical hedging problem has four simultaneous challenges:
(1) discrete rebalancing,
(2) transaction costs,
(3) model misspecification,
(4) complex payoffs.
Black--Scholes delta addresses none of them.
We need a framework that can jointly optimize over all four.
\end{keyresult}

%% ============================================================================
\section{Deep Hedging Framework (Buehler et al., 2019)}
\label{sec:deep-hedging}
%% ============================================================================

The central idea of deep hedging is disarmingly simple:
\textbf{replace the hedging rule with a neural network and train it to minimize
a risk measure of the hedging P\&L.}

\subsection{The Setup}

Consider a discrete-time market with rebalancing dates $t_0, t_1, \ldots, t_N$.
At each time $t_k$, the hedger observes a \emph{state vector}:
\[
    \bm{x}_k = \begin{pmatrix} S_{t_k} \\ \delta_{k-1} \\ T - t_k \\ \sigma_{t_k}^{\text{impl}} \\ \ldots \end{pmatrix}
\]
which may include the current stock price, the current hedge position, time to maturity,
implied volatility, and any other relevant market information.

A \textbf{hedging policy} is a function $\pi_\theta : \bm{x}_k \mapsto \delta_k$
that maps the current state to the number of shares to hold.
Here $\theta$ denotes the parameters of a neural network.

The hedging P\&L for a single path is:
\[
    \PnL(\theta) = \underbrace{C_0}_{\text{premium received}}
    - \underbrace{\max(S_T - K, 0)}_{\text{payoff owed}}
    + \underbrace{\sum_{k=0}^{N-1} \delta_k(\theta) \left(S_{t_{k+1}} - S_{t_k}\right)}_{\text{hedging gains}}
    - \underbrace{\sum_{k=0}^{N-1} c \cdot |\delta_k(\theta) - \delta_{k-1}(\theta)| \cdot S_{t_k}}_{\text{transaction costs}}
\]
where $\delta_k(\theta) = \pi_\theta(\bm{x}_k)$ is the network's output at step $k$.

\subsection{The Optimization Problem}

The network is trained to solve:
\[
    \min_\theta \; \rho\!\left(-\PnL(\theta)\right)
\]
where $\rho$ is a \textbf{convex risk measure} applied to the loss $L = -\PnL$.

\begin{keyidea}
The choice of $\rho$ determines the hedging behavior:
\begin{itemize}
    \item $\rho = \operatorname{Var}$: minimize the variance of the hedging P\&L.
    \item $\rho = \CVaR_\alpha$: minimize the expected loss in the worst $\alpha$-fraction of outcomes.
    \item $\rho = \rho_\lambda^{\text{ent}}$: entropic risk measure with risk aversion $\lambda$,
          i.e., $\rho_\lambda(L) = \frac{1}{\lambda}\ln \EV[e^{\lambda L}]$.
\end{itemize}
The optimization is over all possible hedging strategies simultaneously ---
not step-by-step as in classical dynamic programming.
\end{keyidea}

\subsection{Why This Works: Universal Approximation}

A feedforward neural network with a single hidden layer of sufficient width
can approximate any continuous function on a compact domain (Cybenko, 1989; Hornik, 1991).
The hedging policy $\delta_k = \pi(\bm{x}_k)$ is a continuous function of the state
(at least for vanilla options away from barriers and expiry),
so a sufficiently expressive network can represent it.

The deeper insight is about \textbf{end-to-end optimization}:
rather than decomposing the problem into (a) model calibration,
(b) Greek computation, and (c) cost optimization as separate steps,
deep hedging optimizes the \emph{entire pipeline} jointly through a single loss function.
This eliminates compounding errors from the multi-step classical approach.

\begin{keyresult}
Deep hedging learns the entire hedging policy at once.
The network does not need to know Black--Scholes, compute Greeks, or solve a PDE.
It discovers the optimal strategy directly from the objective and simulated data.
In the zero-cost limit, it \emph{recovers} the Black--Scholes delta.
With costs, it discovers strategies that have no known closed-form solution.
\end{keyresult}

%% ============================================================================
\section{The Training Pipeline}
\label{sec:training-pipeline}
%% ============================================================================

The practical implementation of deep hedging follows a simulation-based training loop.

\subsection{Step 1: Path Simulation}

Generate $M$ sample paths of the underlying price process.
For a GBM model:
\[
    S_{t_{k+1}} = S_{t_k} \exp\!\left[\left(r - \frac{\sigma^2}{2}\right)\Delta t
    + \sigma \sqrt{\Delta t}\; Z_k\right], \quad Z_k \sim \mathcal{N}(0,1)
\]
For more complex dynamics (Heston stochastic volatility, jump-diffusion, rough volatility),
the simulation scheme changes but the training pipeline remains identical.

\begin{keyidea}
The model used for simulation is a \textbf{design choice}, not a constraint.
Deep hedging does not require the simulation model to be ``correct'' ---
it finds the best hedge \emph{within} the simulated environment.
If the simulation model is close to reality, the learned policy will transfer well.
One can also train on an ensemble of models for robustness.
\end{keyidea}

\subsection{Step 2: Forward Pass Through the Network}

For each simulated path $i = 1, \ldots, M$:
\begin{enumerate}
    \item Initialize the position: $\delta_{-1}^{(i)} = 0$ (no initial stock holding).
    \item For each time step $k = 0, 1, \ldots, N-1$:
    \begin{enumerate}
        \item Construct the state vector $\bm{x}_k^{(i)} = (S_{t_k}^{(i)}, \delta_{k-1}^{(i)}, T - t_k, \ldots)$.
        \item Compute $\delta_k^{(i)} = \pi_\theta(\bm{x}_k^{(i)})$.
    \end{enumerate}
    \item Compute the hedging P\&L:
    \[
        \PnL^{(i)}(\theta) = C_0 - \max(S_T^{(i)} - K, 0)
        + \sum_{k=0}^{N-1} \delta_k^{(i)}(S_{t_{k+1}}^{(i)} - S_{t_k}^{(i)})
        - c \sum_{k=0}^{N-1} |\delta_k^{(i)} - \delta_{k-1}^{(i)}| \cdot S_{t_k}^{(i)}
    \]
\end{enumerate}

\subsection{Step 3: Loss Computation}

Compute the risk measure over the batch.
For the variance objective:
\[
    \mathcal{L}_{\text{Var}}(\theta) = \frac{1}{M}\sum_{i=1}^{M}\left(\PnL^{(i)}(\theta)\right)^2
    - \left(\frac{1}{M}\sum_{i=1}^{M}\PnL^{(i)}(\theta)\right)^2
\]

For the CVaR objective at level $\alpha$ (e.g., $\alpha = 0.05$):
\[
    \mathcal{L}_{\CVaR}(\theta) = -\frac{1}{\lfloor \alpha M \rfloor}
    \sum_{j=1}^{\lfloor \alpha M \rfloor} \PnL_{(j)}(\theta)
\]
where $\PnL_{(1)} \leq \PnL_{(2)} \leq \cdots$ are the ordered P\&L values
(the negative sign converts the worst P\&Ls into a quantity we minimize).

In practice, a differentiable approximation of CVaR is used:
\[
    \CVaR_\alpha(L) = \min_{v \in \RV}\left\{v + \frac{1}{\alpha}\EV\!\left[(L - v)^+\right]\right\}
\]
where $L = -\PnL$ is the loss.
This Rockafellar--Uryasev representation is smooth and amenable to gradient-based optimization.

\subsection{Step 4: Backpropagation and Update}

Compute the gradient $\nabla_\theta \mathcal{L}(\theta)$ via backpropagation
through the entire trajectory (all $N$ time steps, all $M$ paths).
Update the parameters with Adam or another stochastic gradient method:
\[
    \theta \leftarrow \theta - \eta \cdot \nabla_\theta \mathcal{L}(\theta)
\]

\subsection{Architecture}

The network architecture is deliberately simple:
\begin{itemize}
    \item \textbf{Input}: state vector of dimension $d$ (typically $d = 3$ to $5$).
    \item \textbf{Hidden layers}: 2--3 fully connected layers with 32--64 units each.
    \item \textbf{Activation}: ReLU or softplus.
    \item \textbf{Output}: a single scalar $\delta_k \in [0, 1]$ (for a long call hedge).
          A sigmoid output layer enforces the bounds.
\end{itemize}

\begin{warningbox}
The same network is applied at every time step --- this is parameter sharing across time,
similar to a recurrent architecture but without hidden state.
The time-to-maturity input $T - t_k$ allows the network to learn time-dependent behavior
(e.g., hedging more aggressively near expiry).
An alternative is to use separate networks per time step,
but shared parameters generalize better with less data.
\end{warningbox}

\begin{refbox}
\begin{itemize}
    \item Buehler, H., Gonon, L., Teichmann, J.\ \& Wood, B.\ (2019).
          ``Deep Hedging.'' \emph{Quantitative Finance}, 19(8), 1271--1291.
          Section~3 details the architecture and training procedure.
    \item Ruf, J.\ \& Wang, W.\ (2020).
          ``Neural Networks for Option Pricing and Hedging: A Literature Review.''
          \emph{Journal of Computational Finance}, 24(1).
          Comprehensive survey of architectures and training tricks.
\end{itemize}
\end{refbox}

%% ============================================================================
\section{Transaction Costs and the Optimal Trade-off}
\label{sec:transaction-costs}
%% ============================================================================

Transaction costs fundamentally alter the hedging problem.
With zero costs, the optimal strategy is to rebalance to the Black--Scholes delta
at every opportunity.
With nonzero costs, the hedger faces a trade-off:
rebalancing reduces hedging error but incurs a direct cost.

\subsection{The No-Trade Zone}

\begin{keyidea}
When trained with transaction costs, the deep hedging network learns to develop
a \textbf{no-trade zone}: a band around the Black--Scholes delta
within which it does not rebalance.
If the current position $\delta_{k-1}$ is within the band, the network outputs
$\delta_k \approx \delta_{k-1}$ (hold steady).
If $\delta_{k-1}$ deviates beyond the band, the network trades toward
(but not all the way to) the Black--Scholes delta.

This behavior emerges entirely from the optimization ---
no one tells the network about no-trade zones.
It rediscovers a result from classical optimal control theory.
\end{keyidea}

Leland (1985) derived an analytical adjustment to the Black--Scholes delta for
proportional transaction costs.
The Leland-adjusted volatility is:
\[
    \hat{\sigma}_L = \sigma\sqrt{1 + \sqrt{\frac{2}{\pi}} \cdot \frac{c}{\sigma\sqrt{\Delta t}}}
\]
and the adjusted delta uses $\hat{\sigma}_L$ in place of $\sigma$.
This gives a wider effective delta range and reduces the rebalancing frequency.

\begin{example}[No-trade zone width]\label{ex:no-trade}
With $\sigma = 20\%$, $\Delta t = 1/365$, and $c = 0.001$ (10 bps):
\[
    \frac{c}{\sigma\sqrt{\Delta t}} = \frac{0.001}{0.20 \times 0.0523} = 0.0957
\]
\[
    \hat{\sigma}_L = 0.20\sqrt{1 + \sqrt{2/\pi} \times 0.0957}
    = 0.20\sqrt{1 + 0.0763} \approx 0.2076
\]

The Leland adjustment widens the effective volatility by about 4\%.
This produces a slight broadening of the delta profile,
which reduces turnover by approximately 8--12\%.

The deep hedging network, when trained under the same conditions,
develops a no-trade zone of comparable width but with \emph{adaptive} boundaries
that depend on the current gamma, time to maturity, and moneyness ---
a level of sophistication that Leland's formula cannot capture.
\end{example}

\subsection{Cost Sensitivity}

As the cost rate $c$ increases:
\begin{enumerate}
    \item The no-trade zone widens.
    \item The hedge ratio deviates more from the Black--Scholes delta.
    \item The hedging P\&L variance increases (accepting more risk to save on costs).
    \item At very high costs, the network may learn to not hedge at all
          (the option premium already compensates for the expected payoff).
\end{enumerate}

\begin{center}
\begin{tabular}{@{}lccc@{}}
\toprule
\textbf{Cost rate} $c$ & \textbf{Avg.\ trades/day} & \textbf{P\&L Std.\ Dev.} & \textbf{Total TC} \\
\midrule
0 bps      & 1.00 & \$0.41 & \$0.00 \\
5 bps      & 0.87 & \$0.44 & \$0.12 \\
10 bps     & 0.71 & \$0.49 & \$0.18 \\
50 bps     & 0.38 & \$0.72 & \$0.31 \\
\bottomrule
\end{tabular}
\end{center}

The table above (representative values from deep hedging simulations) shows the trade-off:
the network smoothly interpolates between frequent rebalancing (low cost) and sparse rebalancing
(high cost), finding the optimal balance that no closed-form solution provides.

\begin{keyresult}
In the zero-cost limit, deep hedging recovers the Black--Scholes delta hedge.
As costs increase, it discovers Leland-type adjustments and extends them
to regimes where no analytical solution exists.
The no-trade zone is not hard-coded --- it \textbf{emerges from optimization}.
\end{keyresult}

%% ============================================================================
\section{Risk Measures as Objectives}
\label{sec:risk-measures}
%% ============================================================================

The choice of risk measure $\rho$ in the deep hedging objective is not merely a
technical detail --- it fundamentally shapes the hedging strategy.

\subsection{Variance Minimization}

The most classical objective:
\[
    \rho_{\text{Var}}(L) = \operatorname{Var}(L) = \EV[L^2] - (\EV[L])^2
\]
where $L = -\PnL$ is the hedging loss.

\textbf{Properties}:
\begin{itemize}
    \item Treats positive and negative deviations symmetrically.
    \item The network penalizes large gains just as much as large losses.
    \item Computationally simple and stable.
    \item Corresponds to quadratic utility: $U(x) = x - \frac{\lambda}{2}x^2$.
\end{itemize}

\begin{warningbox}
Variance penalizes \emph{upside} P\&L deviations equally with \emph{downside}.
A hedging strategy that occasionally produces large profits is penalized just as
severely as one that produces large losses.
For most practitioners, this is undesirable: they care about downside risk, not upside.
\end{warningbox}

\subsection{Conditional Value-at-Risk (CVaR)}

The CVaR at level $\alpha$ (also called Expected Shortfall) is:
\[
    \CVaR_\alpha(L) = \EV\!\left[L \mid L \geq \VaR_\alpha(L)\right]
    = \min_{v \in \RV}\left\{v + \frac{1}{\alpha}\EV\!\left[(L - v)^+\right]\right\}
\]
The second representation (Rockafellar \& Uryasev, 2000) is differentiable
and convenient for optimization.

\textbf{Properties}:
\begin{itemize}
    \item Focuses exclusively on the \emph{worst $\alpha$-fraction} of outcomes.
    \item At $\alpha = 5\%$: ``minimize the average loss in the worst 5\% of scenarios.''
    \item Coherent risk measure (subadditive, monotone, translation-invariant, positive homogeneous).
    \item Produces hedges that are more conservative in the tail.
\end{itemize}

\subsection{Entropic Risk Measure}

The entropic risk measure with risk aversion parameter $\lambda > 0$:
\[
    \rho_\lambda^{\text{ent}}(L) = \frac{1}{\lambda}\ln\EV\!\left[e^{\lambda L}\right]
\]

\textbf{Properties}:
\begin{itemize}
    \item Exponential tilting toward large losses.
    \item As $\lambda \to 0$: $\rho_\lambda^{\text{ent}}(L) \to \EV[L]$ (risk-neutral).
    \item As $\lambda \to \infty$: $\rho_\lambda^{\text{ent}}(L) \to \operatorname{ess\,sup}(L)$ (worst-case).
    \item Corresponds to exponential utility: $U(x) = -e^{-\lambda x}$.
    \item Convex risk measure (but not positively homogeneous).
\end{itemize}

\subsection{How the Objective Shapes the Hedge}

\begin{keyidea}
Different risk measures produce qualitatively different hedging strategies:
\begin{enumerate}
    \item \textbf{Variance}: the network hedges symmetrically, aiming for a tight P\&L distribution
          centered near zero. It may accept some tail risk if doing so reduces overall variance.
    \item \textbf{CVaR}: the network over-hedges slightly, buying extra protection against
          extreme adverse moves. P\&L mean may be slightly negative (paying a ``premium'' for
          tail safety), but the worst outcomes are significantly improved.
    \item \textbf{Entropic}: behavior interpolates between mean and worst-case depending on $\lambda$.
          Higher $\lambda$ produces more conservative hedges, converging toward CVaR-like behavior.
\end{enumerate}
\end{keyidea}

The table below illustrates how the P\&L distribution changes with the objective
(representative values for a European call with $c = 0.1\%$):

\begin{center}
\begin{tabular}{@{}lcccc@{}}
\toprule
\textbf{Objective} & \textbf{Mean P\&L} & \textbf{Std.\ Dev.} & \textbf{5\% CVaR} & \textbf{1\% CVaR} \\
\midrule
Variance    & $-\$0.05$ & \$0.46 & $-\$0.98$ & $-\$1.41$ \\
CVaR (5\%)  & $-\$0.09$ & \$0.52 & $-\$0.79$ & $-\$1.15$ \\
Entropic ($\lambda=2$)  & $-\$0.07$ & \$0.49 & $-\$0.87$ & $-\$1.24$ \\
\bottomrule
\end{tabular}
\end{center}

\begin{remark}
The CVaR-trained hedge has a slightly lower mean and higher variance than the variance-trained hedge,
but its tail risk (5\% CVaR) is substantially better: $-\$0.79$ versus $-\$0.98$,
a 19\% reduction in expected tail loss.
This is the classic risk--return trade-off embedded in the choice of objective.
\end{remark}

%% ============================================================================
\section{Comparison: Deep Hedging vs.\ Delta Hedging}
\label{sec:comparison}
%% ============================================================================

We now present a comprehensive numerical comparison.

\subsection{Setup}

\begin{itemize}
    \item \textbf{Option}: European call, $S_0 = 100$, $K = 100$, $\sigma = 20\%$,
          $r = 0$, $T = 30/365$.
    \item \textbf{Rebalancing}: daily ($N = 30$).
    \item \textbf{Transaction cost}: $c = 0.1\%$ (10 basis points).
    \item \textbf{Paths}: 100{,}000 for training, 50{,}000 for evaluation (out-of-sample).
    \item \textbf{Network}: 3 hidden layers of 64 units, ReLU activation,
          sigmoid output, shared across time steps.
    \item \textbf{Training}: Adam optimizer, learning rate $10^{-3}$, 200 epochs, batch size 1{,}024.
\end{itemize}

\subsection{Results}

\begin{example}[Full comparison]\label{ex:comparison}
The table below reports out-of-sample performance over 50{,}000 paths.

\begin{center}
\begin{tabular}{@{}lccccc@{}}
\toprule
\textbf{Strategy} & \textbf{Mean} & \textbf{Std} & \textbf{5\% CVaR} & \textbf{1\% CVaR} & \textbf{Total TC} \\
\midrule
BS Delta     & $-\$0.18$ & \$0.49 & $-\$1.12$ & $-\$1.58$ & \$0.21 \\
DH (Var)     & $-\$0.06$ & \$0.44 & $-\$0.95$ & $-\$1.37$ & \$0.14 \\
DH (CVaR 5\%)  & $-\$0.11$ & \$0.48 & $-\$0.81$ & $-\$1.12$ & \$0.16 \\
DH (CVaR 1\%)  & $-\$0.14$ & \$0.51 & $-\$0.86$ & $-\$1.02$ & \$0.17 \\
\bottomrule
\end{tabular}
\end{center}

\textbf{Observations}:
\begin{enumerate}
    \item \textbf{Mean P\&L}: BS delta hedging loses \$0.18 on average due to transaction costs
          incurred from daily rebalancing.
          The variance-trained network loses only \$0.06 --- it rebalances less often, saving \$0.07 in costs.
    \item \textbf{Standard deviation}: the DH (Var) network reduces std from \$0.49 to \$0.44 (10\% reduction).
    \item \textbf{5\% CVaR}: the CVaR-trained network achieves $-\$0.81$ vs.\ $-\$1.12$ for BS,
          a \textbf{28\% reduction} in expected tail loss.
    \item \textbf{1\% CVaR}: the DH (CVaR 1\%) network achieves $-\$1.02$ vs.\ $-\$1.58$ for BS,
          a \textbf{35\% reduction}.
    \item \textbf{Transaction costs}: all DH strategies pay less in total costs than BS,
          because they learn to avoid unnecessary rebalancing.
\end{enumerate}
\end{example}

\begin{keyresult}
Deep hedging consistently outperforms Black--Scholes delta hedging in the presence of
transaction costs:
\begin{itemize}
    \item With the variance objective: lower variance, lower mean loss, lower costs.
    \item With the CVaR objective: dramatically better tail performance (15--35\% CVaR reduction),
          at the cost of slightly higher variance.
\end{itemize}
The improvement is most pronounced in the tails,
precisely where risk management matters most.
\end{keyresult}

\subsection{Visualizing the Learned Policy}

The learned hedge ratio $\delta_k(\theta)$ can be plotted against the moneyness $S/K$
at different times to maturity.
Key observations:
\begin{itemize}
    \item \textbf{Zero cost}: the learned curve closely matches the BS delta curve $\Phi(d_1)$.
    \item \textbf{Moderate cost}: the curve flattens (lower gamma sensitivity),
          and the network's output is ``stickier'' --- it changes less between time steps.
    \item \textbf{Near expiry}: the BS delta has a sharp step function near $S = K$.
          The network smooths this transition, avoiding the expensive rapid rebalancing
          that BS would demand.
\end{itemize}

%% ============================================================================
\section{Extensions and Frontier}
\label{sec:extensions}
%% ============================================================================

The basic deep hedging framework generalizes in several important directions.

\subsection{Model-Free (Data-Driven) Hedging}

Instead of simulating from a parametric model (GBM, Heston),
one can train the network directly on historical price paths.
This removes the model specification entirely:
\begin{itemize}
    \item \textbf{Advantage}: no model misspecification risk.
    \item \textbf{Challenge}: limited historical data, especially for rare events.
    \item \textbf{Solution}: augment historical paths with synthetic data from GANs
          or time-series bootstrap methods.
\end{itemize}

\subsection{Rough Volatility Hedging}

Under rough volatility models (e.g., the rough Bergomi model where the volatility process
has Hurst parameter $H \approx 0.1$), the option price is not a Markov function
of the current stock price and time.
Classical delta hedging breaks down entirely.
Deep hedging handles this naturally: the state vector can include past volatility values
or other path-dependent features, and the network learns the correct hedge
\emph{without} requiring Markovianity.

Bayer, Friz \& Gatheral (2016) established the rough volatility paradigm.
Subsequent work has shown that deep hedging under rough volatility outperforms
classical hedges by an even wider margin than in the smooth-volatility case.

\subsection{Multi-Asset Hedging}

When hedging a basket option or a portfolio of derivatives,
the state space grows with the number of assets.
The network input becomes $(S_1, S_2, \ldots, S_d, \delta_1, \ldots, \delta_d, \tau, \ldots)$,
and the output is a vector of hedge ratios.
Deep hedging scales to this setting more gracefully than classical PDE-based methods,
which suffer from the curse of dimensionality.

\subsection{Hedging with Multiple Instruments}

In practice, hedgers use not just the underlying stock but also other options,
futures, or variance swaps.
The deep hedging framework accommodates this by extending the output to
$(\delta_k^{\text{stock}}, \delta_k^{\text{option}}, \ldots)$,
one hedge ratio per instrument.
The network jointly optimizes the allocation across instruments ---
for example, using an option to hedge vega exposure while using the stock for delta.

\subsection{Recent Developments (2021--2025)}

\begin{itemize}
    \item \textbf{Cao, Chen, Hull \& Poulos (2021)}: deep hedging with reinforcement learning,
          comparing policy gradient methods with the direct optimization of Buehler et al.
    \item \textbf{Horvath, Teichmann \& \v{Z}uri\v{c} (2021)}: deep hedging under rough volatility,
          demonstrating that neural hedges outperform classical hedges by a larger margin
          when the volatility surface exhibits rough behavior.
    \item \textbf{Murray, Wood \& Buehler (2022)}: deep hedging for exotic options
          (barrier options, autocallables), showing that the framework extends seamlessly
          to path-dependent payoffs.
    \item \textbf{Gierjatowicz et al.\ (2023)}: robust deep hedging using adversarial training,
          where the network is trained against a worst-case model within an ambiguity set.
    \item \textbf{Horel, Leal \& Ribeiro (2025)}: scaling deep hedging to large derivative portfolios
          using attention-based architectures and transfer learning across maturities.
\end{itemize}

\begin{prereq}
Before implementing deep hedging, ensure you are comfortable with:
\begin{itemize}
    \item Black--Scholes pricing and the Greeks (delta, gamma, vega).
    \item Monte Carlo simulation of asset price paths.
    \item Neural network training with PyTorch or TensorFlow (forward pass, loss, backpropagation).
    \item Convex risk measures: variance, VaR, CVaR (see Project~05 for EVT-based risk measures).
\end{itemize}
\end{prereq}

%% ============================================================================
\section{References}
\label{sec:references}
%% ============================================================================

The following texts provide the mathematical foundations and practical guidance
for deep hedging and its extensions.

\begin{enumerate}[label={[\arabic*]}, leftmargin=2cm]
    \item \textbf{Buehler, H., Gonon, L., Teichmann, J.\ \& Wood, B.}\ (2019).
          ``Deep Hedging.''
          \emph{Quantitative Finance}, 19(8), 1271--1291.
          The foundational paper.
          Introduces the framework, proves convergence,
          and demonstrates results for European and exotic options under transaction costs.

    \item \textbf{Leland, H.E.}\ (1985).
          ``Option Pricing and Replication with Transactions Costs.''
          \emph{The Journal of Finance}, 40(5), 1283--1301.
          Classical treatment of discrete hedging with proportional transaction costs.
          The adjusted volatility formula and no-trade zone concept.

    \item \textbf{Cao, J., Chen, J., Hull, J.\ \& Poulos, Z.}\ (2021).
          ``Deep Hedging of Derivatives Using Reinforcement Learning.''
          \emph{The Journal of Financial Data Science}, 3(1), 10--27.
          Compares deep hedging with reinforcement learning approaches.

    \item \textbf{Horvath, B., Teichmann, J.\ \& \v{Z}uri\v{c}, \v{Z}.}\ (2021).
          ``Deep Hedging under Rough Volatility.''
          \emph{Risks}, 9(7), 138.
          Extends deep hedging to the rough Bergomi model;
          shows neural hedges outperform classical hedges by a wider margin
          in the rough volatility setting.

    \item \textbf{Ruf, J.\ \& Wang, W.}\ (2020).
          ``Neural Networks for Option Pricing and Hedging: A Literature Review.''
          \emph{Journal of Computational Finance}, 24(1).
          Comprehensive survey of neural-network methods in derivatives.

    \item \textbf{Rockafellar, R.T.\ \& Uryasev, S.}\ (2000).
          ``Optimization of Conditional Value-at-Risk.''
          \emph{Journal of Risk}, 2(3), 21--41.
          The CVaR optimization representation used in deep hedging training.

    \item \textbf{Bayer, C., Friz, P.\ \& Gatheral, J.}\ (2016).
          ``Pricing under Rough Volatility.''
          \emph{Quantitative Finance}, 16(6), 887--904.
          Establishes the rough volatility paradigm that motivates
          non-Markovian deep hedging extensions.

    \item \textbf{Gierjatowicz, P.\ et al.}\ (2023).
          ``Robust Deep Hedging.''
          \emph{Quantitative Finance}, 23(7--8), 1079--1098.
          Adversarial training for model-robust hedging strategies.

    \item \textbf{Horel, T., Leal, L.\ \& Ribeiro, R.}\ (2025).
          ``Scaling Deep Hedging to Large Portfolios.''
          \emph{Risk}, January 2025.
          Attention-based architectures and transfer learning for portfolio-level hedging.
\end{enumerate}

\vfill
\hrule
\smallskip
\noindent\textit{Andr\'es Gonz\'alez Ortega --- Risk Analyst Project} \hfill \textit{March 2026}

\end{document}
